\documentclass[10pt,a4paper]{amsart}
\usepackage[margin=19mm]{geometry}
\usepackage{amsmath,amssymb,mathtools}
\usepackage[scr=rsfs,cal=euler]{mathalfa}
\usepackage{xfrac}
\usepackage[unicode,hidelinks,bookmarks=false]{hyperref}
\newcommand{\ie}{i.e.\ }
\newcommand{\cf}{cf.\ }
\newcommand{\resp}{respectively\ }
\newcommand{\Subsection}{\subsection}
\providecommand{\nobreakdash}{\nobreak\hbox{-}}
\setlength{\emergencystretch}{2em}
\multlinegap0pt
\usepackage[shortlabels]{enumitem}
\usepackage{amssymb}
\usepackage{mathtools}
\newtheorem{theorem}{Theorem}
\newtheorem{proposition}{Proposition}
\newcommand{\Eff}{{\rm {Eff}}}
\newcommand{\Exc}{{\rm {Exc}}}
\newcommand{\Mov}{{\rm {Mov}}}
\newcommand{\NE}{{\rm {NE}}}
\newcommand{\Nef}{{\rm {Nef}}}
\newcommand{\bQ}{\mathbb{Q}}
\newcommand{\bR}{\mathbb{R}}
\newcommand{\bc}{{\rm {bc}}}
\newcommand{\mD}{\mathcal{D}}
\newcommand{\mO}{\mathcal{O}}
\title{Variation of cones of divisors in a family of varieties -- Fano type case}
\author{Sung Rak Choi, Zhan Li, Chuyu Zhou}
\date{Source: arXiv:2504.04109v6 (CC BY 4.0)}
\begin{document}
\maketitle

\noindent\emph{Source and licence.} Variation of cones of divisors in a family of varieties -- Fano type case. Version arXiv:2504.04109v6 (CC BY 4.0), distributed by arXiv under the Creative Commons Attribution 4.0 International licence, http://creativecommons.org/licenses/by/4.0/, as recorded in the arXiv licence metadata for this exact version and shown on the abstract page https://arxiv.org/abs/2504.04109v6. Full licence notice: \texttt{licenses/arxiv-2504.04109-CC-BY-4.0.txt}.\par\smallskip

\noindent\textit{Editorial scope.} Counted unit: the article's proposition Eff (source0.tex lines 1282--1288). The article's complete original proof of it is delivered with it (source0.tex lines 1290--1330, one \texttt{proof} environment of 41 lines). No mathematics is written, repaired or re-derived: the statement and the proof are the article's own lines, in the article's own notation, and no retained line is substituted or re-flowed.  Retained uncounted as the source-local context the proof needs: lines 327--338; lines 396--410; lines 1238--1243; lines 1272--1277.  Nothing was omitted from any retained span, so the package carries no omission record.  No retained line contains a citation, so the wrapper carries no bibliography and no bibliography entry.  No retained line contains a figure environment, an image-inclusion command or drawing code, so no image is delivered and the package declares no asset dependency.  Editorial formatting only, in addition: an AMS \texttt{amsart} preamble that restates the article's own declarations and macros used by the retained lines, namely the environments \texttt{theorem}, \texttt{proposition} on separate counters, together with the standard packages those lines need.  The retained environments print in this document as Theorem 1, Proposition 1 in order of appearance, whereas the article prints the same items with the numbering of its own full document.  The complete native source of the article as distributed in the arXiv e-print for this exact version is bundled unchanged as \texttt{sources/176031/source0.tex}.
\begin{theorem}\label{thm: FT is constructible}
Let $f: X \to T$ be a projective, surjective morphism between varieties such that $f_*\mO_X=\mO_T$. Suppose that there exists a Zariski dense subset $S\subset T$ such that for any $s\in S$, the fiber $X_s$ is a Fano type variety. Then, after shrinking $T$, there exists a divisor $D$ on $X$ such that $(X, D)$ has lc singularities and
\[
K_X+D \sim_\bQ 0/T.
\] Furthermore, if for any $s\in S$, the fiber $X_s$ falls into one of the following cases:
\begin{enumerate}[label=(\roman*)]
\item $X_s$ is a klt weak Fano variety for any $s\in S$,
\item $X_s$ is of Fano type for any $s\in S$ and $S$ consists of very general points,
\item $X_s$ is of Fano type for any $s\in S$ and $\{X_s \mid s\in S\}$ admits bounded klt complements,
\end{enumerate}
then, after shrinking $T$, we have that $X$ is of Fano type over $T$.
\end{theorem}

\begin{theorem}\label{mainthm: mmp}
Let $f: X \to T$ be a projective fibration. Suppose that there exists a Zariski dense subset $S\subset T$ such that for any $s\in S$, the fiber $X_s$ satisfies one of the cases in Theorem \ref{thm: FT is constructible}. Then, after a generically finite base change of $T$, there exists a non-empty Zariski open subset $T_0 \subset T$ such that any Zariski open subset $U \subset T_0$ satisfies the following properties.
\begin{enumerate}
\item Suppose that $X'/U \in \bc(X_U/U)$. If $g: X' \dashrightarrow Y/U$ is a birational contraction with $Y$ a $\bQ$-factorial variety, then for any $t \in U$, the induced map $g_t: X'_t \dashrightarrow Y_t$ is again a birational contraction. Moreover, for any $\bR$-Cartier divisor $\mD$ on $X'$, we have
\[
(g_t)_*(\mD|_{X'_t}) \sim_{\bR} (g_*\mD)|_{Y_t}, \quad t \in U.
\]
\item For any $\bR$-divisor $\mD\in N^1(X_U/U)$, a sequence of $\mD$-MMP$/U$ induces a sequence of $\mD|_{X_t}$-MMP of the same type for each $t\in U$. 
\item Conversely, any $\bR$-divisor $D\in N^1(X_t)$ with $t\in U$, a sequence of $D$-MMP on $X_t$ is induced by a sequence of $\mD$-MMP$/U$ on $X_U/U$ of the same type. Moreover, we can choose $\mD$ such that $\mD|_{X_t}\sim_\bR D$. 
\item If $Y/U\in \bc(X_{U}/U)$, then the natural maps
\[
N^1(Y/U) \to N^1(Y_t), \quad \Nef(Y/U) \to \Nef(Y_t), \quad \NE(Y_t) \to \NE(Y/U)
\] are isomorphisms for any $t\in U$.
\end{enumerate}
\end{theorem} 

\begin{proposition}\label{prop: mov}
Let $f: X \to T$ be a projective fibration. Suppose that there exists a subset $S\subset T$ such that the fibers $X_s$ for $s\in S$ satisfy one of the cases in Theorem \ref{thm: FT is constructible}. Then, up to a generically finite base change of $T$, there exists a Zariski open subset $T_0\subset T$, such that for any $U\subset T_0$, we can identify the Mori chamber decompositions of $\Mov(X_{T_0}/T_0), \Mov(X_{U}/U)$ and $ \Mov(X_t), t\in U$ under the natural restriction maps. In particular, we have isomorphisms among movable cones 
\[
\Mov(X_{T_0}/T_0) \to \Mov(X_U/U) \to \Mov(X_t), t\in U
\] under the natural restriction maps.
\end{proposition}

\begin{equation}\label{eq: chamber of eff}
\Eff(X/T)
=
\bigcup_{g\colon X \dashrightarrow Y/T}
\left(g^*\Nef(Y/T) + \Exc(g)\right),
\end{equation}

\begin{proposition}\label{prop: Mori chamber for Eff}
Under the same assumptions as in Proposition \ref{prop: mov}, after a
generically finite base change of $T$, there exists a non-empty Zariski open
subset $T_0\subset T$ such that, for any open subset $U\subset T_0$ and any
$t\in U$, the chamber decompositions in \eqref{eq: chamber of eff} of $
\Eff(X_{T_0}/T_0), \Eff(X_U/U)$ and $\Eff(X_t)$ are identified under the natural restriction maps.
\end{proposition}

\begin{proof}
After replacing $T$ by a generically finite base change and possibly shrinking
it, we may assume that Theorem \ref{mainthm: mmp} holds. By Theorem
\ref{thm: FT is constructible}, after shrinking $T$ further, we may assume that
$X/T$ is of Fano type. By Theorem \ref{mainthm: mmp} (2) and (3), an MMP$/T$ on $X$ restricts to an
MMP of the same type on each fiber $X_t$, and conversely every MMP on $X_t$ is
induced by an MMP$/T$. By Theorem
\ref{mainthm: mmp} (4), for any $Y/U\in \bc(X_U/U)$ and any $t\in U$, the
natural maps
\[
N^1(Y/U) \to N^1(Y_t), \quad
\Nef(Y/U) \to \Nef(Y_t)
\]
are isomorphisms.

It remains to compare the exceptional parts of the chambers. Shrinking $T_0$
further, we may assume that, for every birational contraction
$g\colon X_{T_0}\dashrightarrow Y/T_0$, every $g$-exceptional prime divisor
dominates $T_0$, and its restriction to each fiber is a $g_t$-exceptional
prime divisor. Conversely, by the compatibility of the MMP with fibers
described above, every $g_t$-exceptional divisor arises in this way. Hence the
cones generated by the exceptional divisors are identified under the natural
restriction maps.

Therefore, for any open subset $U\subset T_0$ and any $t\in U$, the
decompositions
\[
\Eff(X_U/U)
=
\bigcup_{g\colon X_U\dashrightarrow Y/U}
\left(g^*\Nef(Y/U)+\Exc(g)\right)
\]
and
\[
\Eff(X_t)
=
\bigcup_{g_t\colon X_t\dashrightarrow Y_t}
\left(g_t^*\Nef(Y_t)+\Exc(g_t)\right)
\]
have the same chambers under the above identifications.
\end{proof}

\end{document}
